\section{What counts as treatment in language data}
\label{sec:treatments}

In a policy evaluation the treatment is usually a known act on a known date: a law takes effect, a price changes, a program opens. Language data rarely offers that. A change in usage may be pushed by a circular, a dictionary, an administrative endorsement, a first enforcement, a media house revising its style sheet, or a social movement that made all of these possible. Which act counts as the treatment isn't given by the corpus; it fixes the estimand.

The shock need not be a policy. It can be a new medium with constraints of its own, a contact or demographic event, a censorship regime imposed or lifted, or a non-linguistic event that drags vocabulary behind it. These can be cleaner than decrees because their timing is less likely to be chosen by the speech community. The policy case this paper uses is the demanding one: a decree's timing may be bound up with the trend it's meant to explain.

The feminization case shows the ambiguity. Candidate treatments include an endorsement, a feminized-titles dictionary \citep{Moreau1991}, a circular, and the date editors actually began applying the forms. France's 1986 circular \citep{france1986circ} urged feminine titles and went largely unenforced, by the government's own later admission, until a second circular in 1998 \citep{france1998circ,BurnettBonami2019}. Those two acts share their text and differ in force, which makes 1986 useful as a near-null policy: a case that should show little or nothing, not a failed treatment.

Official dates can also lag the changes they name. An endorsement often ratifies a shift already under way rather than starting it, a question of actuation \citep{BurnettBonami2019}. Feminization is a socially structured diffusion, an S-curve spreading token by token and register by register \citep{weinreich1968,tagliamontedarcy2007}.

Two varieties can both be rising and still violate parallel trends if one is in the steep middle of its curve and the other is leveling off. A shock identifies an effect only when it appears as a level shift or acceleration on top of comparable untreated trajectories; otherwise \enquote{not identified} is a property of the design, not a disappointment after estimation.

Treatments in this domain are rarely binary either. Endorsements vary in scope and force, so two polities labelled treated may have received different doses. Where dose is measurable, intensity is more honest than an on-off indicator; where it isn't, the binary coding is an assumption to report.

The boundary also leaks: a policy in one polity reaches readers in another through shared media and the channels opened by a common language. The first question of a corpus DiD is this: which act, at which date, with what force and reach, is being called the cause. The estimand, comparison, and diagnostics inherit that choice.
